\section{Методы оптимизации}
\label{sec:methods}

Обучение машинно-обучаемого потенциала будем рассматривать как минимизацию функции потерь по параметрам модели. Пусть \(\theta\) обозначает вектор всех обучаемых параметров потенциала, а \(\mathcal L(\theta)\) --- функционал ошибки по энергиям и силам. Тогда задача обучения записывается как
\begin{equation}
    \theta^\ast
    =
    \arg\min_{\theta}
    \mathcal L(\theta).
\end{equation}
В работе намеренно сравниваются методы с разным объёмом информации о функции потерь. Методы первого порядка опираются только на градиент. Квазиньютоновские методы дополнительно накапливают приближение кривизны. Отдельно рассматривается Muon: в нём матричное направление обновления нормализуется с помощью ортогонализации.

\subsection{Градиентный спуск и стохастический градиентный спуск}

Самый простой метод первого порядка -- градиентный спуск. На каждой итерации параметры обновляются в направлении антиградиента:
\begin{equation}
    \theta_{t+1}
    =
    \theta_t
    -
    \eta_t
    \nabla \mathcal L(\theta_t),
\end{equation}
где \(\eta_t>0\) --- шаг обучения. Такой метод использует полный градиент по всей обучающей выборке, поэтому при фиксированных параметрах направление шага детерминировано. Недостаток состоит в стоимости одной итерации: она растёт вместе с числом конфигураций и числом атомов.

Стохастический градиентный спуск использует градиент не по всей выборке, а по мини-пакету \(\mathcal B_t\):
\begin{equation}
    g_t
    =
    \nabla
    \mathcal L_{\mathcal B_t}(\theta_t),
    \qquad
    \theta_{t+1}
    =
    \theta_t
    -
    \eta_t g_t.
\end{equation}
Такой шаг дешевле, но \(g_t\) является шумной оценкой полного градиента. Поэтому SGD вводит дополнительный компромисс: стоимость итерации уменьшается, а дисперсия направления шага растёт. Для обучения межатомных потенциалов это особенно заметно, поскольку вклад сил даёт большое число слагаемых в функции потерь, а масштабы градиентов по разным параметрам могут существенно отличаться.

На практике результат SGD сильно зависит от learning rate, размера mini-batch и ограничения нормы градиента. Слишком большой шаг может увести параметры в область больших значений loss, а слишком малый делает обучение медленным и часто приводит к затяжному плато.

\subsection{Momentum}

Momentum расширяет SGD за счёт накопления направления движения. Вместо использования только текущего градиента вводится вспомогательная переменная \(v_t\):
\begin{equation}
    v_{t+1}
    =
    \beta v_t
    +
    g_t,
\end{equation}
\begin{equation}
    \theta_{t+1}
    =
    \theta_t
    -
    \eta_t v_{t+1},
\end{equation}
где \(\beta\in[0,1)\) --- коэффициент инерции. Если последовательные градиенты направлены примерно одинаково, величина \(v_t\) накапливается, и движение в этом направлении ускоряется. Если градиенты колеблются, усреднение частично подавляет зигзагообразное движение.

Momentum полезен на плохо обусловленных поверхностях функции потерь. В таких задачах кривизна в разных направлениях может сильно отличаться: в одних направлениях функция меняется резко, в других --- медленно. Обычный SGD в этой ситуации делает много поперечных колебаний. Momentum сглаживает направление шага и может ускорять движение вдоль вытянутых долин функции потерь. При этом метод всё равно сохраняет один общий learning rate и не снимает полностью проблему разных масштабов координат.

\subsection{Adam}

Adam использует адаптивный шаг. В отличие от SGD и Momentum, он хранит две экспоненциально сглаженные статистики: первый момент градиента и второй момент покоординатного квадрата градиента. Для mini-batch градиента \(g_t\) обновления имеют вид
\begin{equation}
    m_t
    =
    \beta_1 m_{t-1}
    +
    (1-\beta_1)g_t,
\end{equation}
\begin{equation}
    v_t
    =
    \beta_2 v_{t-1}
    +
    (1-\beta_2)(g_t\odot g_t),
\end{equation}
где \(\odot\) обозначает покоординатное произведение. Поскольку в начале обучения \(m_t\) и \(v_t\) смещены к нулю, используются bias-correction оценки:
\begin{equation}
    \widehat m_t
    =
    \frac{m_t}{1-\beta_1^t},
    \qquad
    \widehat v_t
    =
    \frac{v_t}{1-\beta_2^t}.
\end{equation}
После этого параметры обновляются по правилу
\begin{equation}
    \theta_{t+1}
    =
    \theta_t
    -
    \eta_t
    \frac{\widehat m_t}
    {\sqrt{\widehat v_t}+\varepsilon}.
\end{equation}
Первый момент \(m_t\) задаёт сглаженное направление, близкое по смыслу к Momentum. Второй момент \(v_t\) нормирует шаг по координатам: если по некоторой координате градиенты систематически велики, эффективный шаг уменьшается; если градиенты малы, относительный шаг увеличивается. За счёт этого Adam частично компенсирует различие масштабов параметров и обычно устойчивее SGD при шумных mini-batch градиентах \cite{kingma2015adam}. В экспериментальной части Adam используется как основной mini-batch baseline.

\subsection{Ограничение нормы градиента}

Для mini-batch методов дополнительно применялось ограничение нормы градиента. Пусть \(C>0\) --- заданный порог clipping. Тогда вместо исходного градиента \(g_t\) используется
\begin{equation}
    g_t^{\mathrm{clip}}
    =
    g_t
    \cdot
    \min
    \left(
        1,
        \frac{C}{\|g_t\|}
    \right).
\end{equation}
Если \(\|g_t\|\leq C\), градиент не изменяется. Если \(\|g_t\|>C\), он масштабируется до нормы \(C\). Направление при этом сохраняется, а максимальный размер обновления ограничивается:
\begin{equation}
    \|\Delta \theta_t\|
    \lesssim
    \eta_t C.
\end{equation}
Для обучения потенциалов такой приём важен из-за резких изменений градиента, связанных с вкладом сил и случайным составом mini-batch. Слишком малое значение \(C\) делает обучение чрезмерно осторожным, а слишком большое почти не влияет на шаг.

\subsection{BFGS}

Методы первого порядка используют только локальный градиент. На плохо обусловленных задачах этой информации часто недостаточно: один и тот же learning rate должен одновременно подходить и для направлений с большой кривизной, и для почти плоских направлений. Квазиньютоновские методы решают эту проблему иначе: они строят приближение кривизны функции потерь и используют его при выборе направления шага.

В методе BFGS направление обновления задаётся формулой
\begin{equation}
    p_k
    =
    -
    H_k
    \nabla \mathcal L(\theta_k),
\end{equation}
где \(H_k\) --- приближение обратного гессиана. Затем выполняется шаг
\begin{equation}
    \theta_{k+1}
    =
    \theta_k
    +
    \alpha_k p_k,
\end{equation}
где \(\alpha_k\) выбирается процедурой line search или заданным правилом выбора шага.

Для обновления приближения обратного гессиана вводятся величины
\begin{equation}
    s_k
    =
    \theta_{k+1}-\theta_k,
    \qquad
    y_k
    =
    \nabla \mathcal L(\theta_{k+1})
    -
    \nabla \mathcal L(\theta_k).
\end{equation}
Обновление BFGS имеет вид
\begin{equation}
    H_{k+1}
    =
    \left(
        I
        -
        \rho_k s_k y_k^\top
    \right)
    H_k
    \left(
        I
        -
        \rho_k y_k s_k^\top
    \right)
    +
    \rho_k s_k s_k^\top,
\end{equation}
где
\begin{equation}
    \rho_k
    =
    \frac{1}{y_k^\top s_k}.
\end{equation}
BFGS не требует явного вычисления полного гессиана, но постепенно накапливает информацию о локальной кривизне через изменения градиента. Поэтому метод хорошо подходит для гладких задач с точным градиентом \cite{nocedal2006numerical}. В обучении MTP и ETN это важно: функция потерь по энергиям и силам гладко зависит от параметров потенциала, а full-batch градиент может быть вычислен средствами MLIP-4.

\subsection{L-BFGS}

Главное ограничение полного BFGS связано с хранением матрицы \(H_k\in\mathbb R^{n\times n}\), где \(n\) --- число параметров модели. Затраты памяти составляют \(O(n^2)\), что для больших моделей быстро становится существенным ограничением.

L-BFGS -- вариант BFGS с ограниченной памятью. Вместо полной матрицы \(H_k\) метод хранит только последние \(m\) пар
\begin{equation}
    (s_{k-m},y_{k-m}),\ldots,(s_{k-1},y_{k-1}).
\end{equation}
Действие приближения обратного гессиана на градиент вычисляется не через явное умножение на матрицу, а через двухпроходную рекурсию по сохранённым векторам. Поэтому память уменьшается до \(O(mn)\), где \(m\ll n\) \cite{nocedal2006numerical}.

В работе L-BFGS рассматривался в комбинированной схеме Adam + L-BFGS. Идея состоит в разделении этапов обучения: Adam сначала выводит параметры в область меньших значений функции потерь, после чего L-BFGS выполняет более точную дооптимизацию с учётом локальной кривизны. Такой подход проверялся как компромисс между устойчивостью адаптивных mini-batch методов и точностью квазиньютоновской оптимизации.

\subsection{Muon}

Muon рассматривался как дополнительный метод для параметров, которые можно представить в матричной форме. В отличие от Adam с покоординатной нормализацией шага, Muon работает с двумерными параметрами как с матрицами и нормализует само матричное направление обновления. В исходной постановке метод применяется к матричным параметрам скрытых слоёв нейронных сетей, а скалярные, векторные, входные и выходные параметры рекомендуется оптимизировать стандартным методом, например AdamW~\cite{jordan2024muon}. Связанные интерпретации матричной нормализации обновлений обсуждаются в работе~\cite{bernstein2024oldoptimizer}.

Базовый шаг Muon начинается с momentum-обновления. Пусть \(G_t\) --- градиент по матричному параметру \(W_t\in\mathbb R^{n\times m}\). Тогда сначала строится momentum-направление
\begin{equation}
    M_t
    =
    \beta M_{t-1}
    +
    G_t.
\end{equation}
Далее это матричное направление ортогонализуется:
\begin{equation}
    \widetilde M_t
    =
    \operatorname{Ortho}(M_t).
\end{equation}
После этого параметр обновляется по правилу
\begin{equation}
    W_{t+1}
    =
    W_t
    -
    \eta_t
    \widetilde M_t.
\end{equation}

Оператор \(\operatorname{Ortho}\) заменяет матрицу обновления на ближайшую полуортогональную матрицу:
\begin{equation}
    \operatorname{Ortho}(G)
    =
    \arg\min_{O}
    \left\{
        \|O-G\|_F:
        O^\top O=I
        \ \text{или}\
        OO^\top=I
    \right\}.
\end{equation}
Если сингулярное разложение матрицы \(G\) имеет вид
\begin{equation}
    G
    =
    U S V^\top,
\end{equation}
то ортогонализованное направление равно
\begin{equation}
    \operatorname{Ortho}(G)
    =
    U V^\top.
\end{equation}
Тем самым Muon сохраняет сингулярные направления обновления, но выравнивает его сингулярные значения. Это уменьшает доминирование нескольких крупных направлений и делает менее выраженные направления матричного обновления более заметными.

Прямое вычисление SVD на каждом шаге слишком дорого, поэтому в Muon ортогонализация приближённо выполняется через итерацию Newton--Schulz. Сначала матрица нормируется:
\begin{equation}
    X_0
    =
    \frac{G}{\|G\|_F+\varepsilon}.
\end{equation}
Затем выполняется несколько итераций вида
\begin{equation}
    A_t
    =
    X_t X_t^\top,
\end{equation}
\begin{equation}
    X_{t+1}
    =
    aX_t
    +
    \left(
        bA_t
        +
        cA_t^2
    \right)
    X_t,
\end{equation}
где в исходной реализации используются коэффициенты
\begin{equation}
    a=3.4445,
    \qquad
    b=-4.7750,
    \qquad
    c=2.0315.
\end{equation}
Если \(G=USV^\top\), то такая итерация действует на сингулярные значения матрицы через полином
\begin{equation}
    \varphi(x)
    =
    ax
    +
    bx^3
    +
    cx^5.
\end{equation}
После нескольких шагов сингулярные значения приближаются к единице, а матрица обновления приближается к \(UV^\top\). Поэтому Muon можно интерпретировать как momentum-метод, в котором матричное направление дополнительно приводится к полуортогональному виду.

В экспериментах Muon проверялся в двух вариантах: \texttt{Muon pad sqrt} и \texttt{Muon factorization}. Здесь важна не теоретическая демонстрация преимущества метода, а его практическое сравнение с SGD, Momentum, Adam и BFGS при обучении MTP и ETN. Результаты показывают, что качество Muon зависит от модели, датасета и метрики: на AgPd метод не превосходит Adam по validation EPA RMSE и уступает BFGS-вариантам, но для ETN даёт меньшую validation forces RMSE, чем Adam; на MoNbTaVW для ETN оба варианта Muon существенно лучше Adam, SGD и Momentum по validation EPA RMSE и дают меньшую energy error, чем BFGS-варианты, хотя по forces RMSE и времени обучения всё ещё уступают full-batch методам.
